% Generated draft: requires mathematical and source review.
\documentclass[11pt]{article}
\usepackage[margin=25mm]{geometry}
\usepackage{fontspec}
\setmainfont{Latin Modern Roman}
\usepackage{amsmath,amssymb,mathtools,unicode-math}
\setmathfont{Latin Modern Math}
\usepackage{enumitem,needspace,xurl,hyperref}
\setlength{\parindent}{0pt}
\setlength{\parskip}{5pt}
\setlength{\emergencystretch}{3em}
\newcommand{\N}{\mathbb{N}}
\newcommand{\Z}{\mathbb{Z}}
\newcommand{\Q}{\mathbb{Q}}
\newcommand{\R}{\mathbb{R}}
\newcommand{\C}{\mathbb{C}}
\newcommand{\F}{\mathbb{F}}
\newcommand{\A}{\mathbb{A}}
\newcommand{\PP}{\mathbb{P}}
\newcommand{\E}{\mathbb{E}}
\newcommand{\eps}{\varepsilon}
\newcommand{\dd}{\,\mathrm{d}}
\begin{document}
\section*{UnsolvedMath ID: 30000537}
\textbf{Source problem code:} catalogue entry
\subsection{Definitions and mathematical statement}
In the usual boundary-value setting, let $(M^{2n},\omega)$ be a symplectic manifold, let $L\subset M$ be a Lagrangian submanifold, and let $J$ be an almost complex structure tamed by $\omega$. A $J$-holomorphic disk is a map $u:(D,\partial D)\to(M,L)$ satisfying
\[
\dd u\circ j=J\circ\dd u.
\]
For a relative class $\beta\in\pi_2(M,L)$ and $r$ boundary and $s$ interior marked points, write $\mathcal M_{r,s}(\beta;J)$ for the corresponding moduli space modulo disk reparametrization, with evaluation maps at the marked points. In the standard index convention its expected real dimension is $n+\mu_L(\beta)+r+2s-3$. Incidence conditions require evaluations to lie on specified cycles in $L$ or $M$; isolated constrained configurations require the total real codimension of those conditions to equal the expected dimension.

The source asks for a construction of an invariant count of these constrained disks that remains meaningful in the presence of bubbling. Disk and sphere bubbles can occur in compactifications and in parameterized families. The requested cluster-moduli-space construction, or a suitable alternative, must account coherently for these degenerations and transversality issues so that the resulting count has the intended invariance. A naive signed count of smooth isolated disks alone is not asserted to have this property.

The extract does not provide $M,L$, the disk classes, incidence cycles, compactness and orientation hypotheses, coefficient ring, allowable auxiliary or bounding data, or the equivalence relation under which the count is to be invariant. The customary moduli-space notation above makes the counting target explicit, without claiming that arbitrary Lagrangians admit an unconditional count. Necessary correction terms, obstruction conditions and any wall-crossing dependence belong to the missing construction; they are not assumed away.
\subsection{Short English statement}
Construct counts of $J$-holomorphic disks satisfying cycle constraints that account for bubbling and are invariant in the intended sense. Cluster spaces or another theory may be used, but the geometric hypotheses, correction data and precise invariance are not specified in the extract.
\subsection{Sources}
\begin{itemize}
\item Supplied report: \emph{Geometric Group Theory, Hyperbolic Dynamics and Symplectic Geometry} (2007), Oberwolfach Reports, DOI 10.4171/owr/2006/33, pp. 1991--2058; proposed year 2006. \url{https://doi.org/10.4171/owr/2006/33}.
\item Supplied reference labelled ``Cluster-homology framework'': \url{https://arxiv.org/abs/math/0404344}.
\item Local UnsolvedMath registry, permanent ID 30000537, catalogue entry; created 25 August 2026, updated 16 September 2026. It labels the statement corrected\_verified following editorial cleanup. Its supplied literature triage of 21 August 2026 labels it partially solved: cluster and virtual-perturbation approaches and important monotone cases are reported, while the general transversality and simultaneous construction are described as incomplete. This dated assessment requires independent verification.
\end{itemize}
\end{document}
